\centerline{\bf \Huge Текстовые задачи II}

\problem{\textit{\textbf{1}}} Найдите отношение двух чисел, если известно, что разность первого числа и $10 \%$ второго числа составляет $50 \%$ суммы второго числа и $50 \%$ первого.

\problem{\textit{\textbf{2}}} Свежие грибы содержат $90 \%$ влаги, сушеные - $12 \%$. Сколько сушеных грибов получится из 10 кг свежих?

\problem{\textit{\textbf{3}}} Антикварный магазин, купив два предмета за 225 р., продал их, получив $40 \%$ прибыли. За какую цену был куплен магазином каждый предмет, если при продаже первого предмета было получено $25 \%$ прибыли, а второго - $50 \%$ ?

\problem{\textit{\textbf{4}}} Гипотенуза прямоугольного треугольника равна $3 \sqrt{ 5}$ м. Найдите катеты, если известно, что после того, как один из них увеличить на $133 \dfrac{1}{3} \%$, а другой - на $16 \dfrac{2}{3} \%$, сумма их длин станет равной 14 м.

\problem{\textit{\textbf{5}}} Каково минимально возможное число учеников в выпускном классе средней школы, если известно, что процент неуспевающих учеников в классе заключен в пределах от $2,5 \%$ до $2,9\%$?

\problem{\textit{\textbf{6}}} Найдите двузначное число, зная, что число его единиц на 2 больше числа десятков, а произведение искомого числа на сумму его цифр равно 280 .



\problem{\textit{\textbf{7}}} Четверо рабочих обрабатывают детали с постоянной производительностью. Если первый будет работать 2 ч, второй 4 ч и четвертый - 6 ч, то вместе они обработают 260 деталей. Если второй и четвертый будут работать по 6 , а третий - 2 ч, то будет обработано 270 деталей. Если второй и четвертый будут работать по 1 ч, то они успеют обработать 40 деталей. Сколько деталей будет обработано, если первый, третий и четвертый рабочий будут работать по 1 ч?

\problem{\textit{\textbf{8}}} Смешали $10 \%$-ный и $25 \%$-ный растворы соли и получили 3 кг $20\%$-ного раствора. Какое количество каждого раствора в килограммах было использовано?

\problem{\textit{\textbf{9}}} Допуская, что стрелки часов движутся без скачков, определите, через какое время после того, как часы показывали 4:00, минутная стрелка догонит часовую.

\problem{\textit{\textbf{10}}} Имеются два сплава золота и серебра. В одном сплаве количество этих металлов находится в отношении $2: 3$, а в другом - в отношении $3: 7$. Сколько нужно взять каждого сплава, чтобы получить 8 кг нового сплава, в котором золото и серебро были бы в отношении $5: 11$ ?

\newpage

\hspace{3mm}

\problem{\textit{\textbf{11}}} Мальчик сбежал вниз по движущемуся эскалатору и насчитал 30 ступенек. Затем он пробежал вверх по тому же эскалатору с той же скоростью относительно эскалатора и насчитал 150 ступенек. Сколько ступенек он насчитал бы, спустившись по неподвижному эскалатору?

\problem{\textit{\textbf{12}}} Из пункта $A$ в пункт $B$, расположенный в 24 км от $A$, одновременно отправились велосипедист и пешеход. Велосипедист прибыл в пункт $B$ на 4 ч раньше пешехода. Известно, что если бы велосипедист ехал с меньшей на 4 км/ч скоростью, то на путь из $A$ в $B$ он затратил бы вдвое меньше времени, чем пешеход. Найдите скорость пешехода.

\problem{\textit{\textbf{13}}} От пристани $A$ отправились одновременно вниз по течению реки катер и плот. Катер спустился вниз по течению на 96 км, затем повернул обратно и вернулся в $A$ через 14ч. Найдите скорость катера в стоячей воде и скорость течения реки, если известно, что катер встретил плот на обратном пути на расстоянии 24 км от $A$.

\problem{\textit{\textbf{14}}} Два поезда отправляются навстречу друг другу из городов $A$ и $B$. Если поезд из города $A$ отправится на 1,5 ч раньше, чем поезд из города $B$, то они встретятся на середине пути. Если оба поезда выйдут одновременно, то через 6 ч они еще не встретятся, а расстояние между ними составит десятую часть первоначального. За сколько часов может проехать каждый поезд расстояние между $A$ и $B$ ?

\problem{\textit{\textbf{15}}} Велосипедист отправляется с некоторой скоростью из пункта $A$ в пункт $B$, отстоящий от $A$ на расстоянии 60 км. Прибыв в $B$, он сразу же выезжает обратно с той же скоростью, но через 1 ч после выезда из $B$ делает остановку на 20 мин, после чего продолжает путь, увеличив скорость на 4 км/ч. В каких пределах заключена скорость велосипедиста, если известно, что на обратный путь от $B$ до $A$ он затратил времени не более, чем на путь от $A$ до $B$.

\problem{\textit{\textbf{16}}} От пристани $A$ вниз по течению реки одновременно отплыли теплоход и плот. Теплоход, доплыв до пристани $B$, расположенной в 324 км от пристани $A$, простоял там 18ч и отправился обратно в $A$. В тот момент, когда он находился в 180 км от $A$, второй теплоход, отплывший из $A$ на 40 ч позднее первого, нагнал плот, успевший к этому времени проплыть 144 км. Считая, что скорость течения реки постоянна, скорость плота равна скорости течения реки, а скорости теплоходов в стоячей воде постоянны и равны между собой, определите скорости теплоходов и скорость течения реки.